% ====================================================================
% Section 4 (v2): The ablation - why two terms and not eight.
% Faithful to data/ablation_raw.json (cap 2,2) and ablation_24_raw.json (cap 2,4).
% Numbers are the change in losses when each layer is REMOVED from the full model.
% ====================================================================
\section{Ablation: Why Two Terms}
\label{sec:ablation}

Section~\ref{sec:core} asserted that only two mechanisms survive scrutiny. This
section is the evidence. EMMET did not begin parsimonious; it began as an
exercise in taking the particle metaphor seriously, and accreted eight physical
layers. We removed them one at a time and measured what each was worth.

\subsection{The candidate layers}

Beyond the gradient (which is the irreducible substrate---without it there is no
router) the full model carried four added mechanisms:

\begin{itemize}
\item \textbf{Newton-III scar} (loss-reaction memory, Section~\ref{sec:core}).
\item \textbf{Effective mass}: a packet's influence on the potential grew with
the congestion it had already traversed, modelling inertia.
\item \textbf{Thermostat}: a global network ``temperature'' that scaled the
congestion weight $\beta$ as mean utilisation rose.
\item \textbf{Temporal ripple}: loss memory deposited not as a single
increment but diffused over several subsequent steps, like a wave.
\end{itemize}

\subsection{Protocol}

We evaluated the full model and four leave-one-out variants (each removing a
single layer), plus a bare core (gradient only, all four layers removed),
against the trivial-but-stateless DRILL baseline on the G\'EANT backbone under
trunk-link failure. Two load regimes were used---moderate (capacity drawn from
$[2,4]$) and saturated ($[2,2]$)---across the failed-link ensemble with 30
seeds each. The quantity of interest is the \emph{change in total losses when a
layer is removed} from the full model: a positive change means the layer was
doing useful work; a negative or zero change means it was inert or harmful.

\subsection{Result: only the scar earns its keep}

Table~\ref{tab:ablation} reports the change in losses on removal of each layer,
in both regimes.

\begin{table}[t]
\centering
\caption{Change in total losses when each layer is removed from the full model
(positive = the layer was helping; negative = it was hurting). G\'EANT, trunk
failure, 30 seeds. The Newton-III scar is the only layer whose removal clearly
hurts; the ripple's removal slightly \emph{helps}.}
\label{tab:ablation}
\small
\begin{tabular}{lcc}
\toprule
Layer removed & moderate load & saturated load \\
\midrule
Temporal ripple   & $-0.4\%$ & $-0.5\%$ \\
Thermostat        & $\phantom{+}0.0\%$ & $\phantom{+}0.0\%$ \\
Effective mass    & $+4.2\%$ & $+1.6\%$ \\
Newton-III scar   & $+16.7\%$ & $+6.4\%$ \\
\midrule
All four (bare core) & $+38.7\%$ & $+9.8\%$ \\
\bottomrule
\end{tabular}
\end{table}

The reading is blunt:

\begin{itemize}
\item The \textbf{Newton-III scar} is the one added layer that earns its keep:
removing it raises losses by $6.4$--$16.7\%$. It is doing the real work.
\item \textbf{Effective mass} contributes marginally ($1.6$--$4.2\%$).
\item The \textbf{thermostat} has \emph{no measurable effect} (0.0\% in both
regimes): it is pure decoration.
\item The \textbf{temporal ripple} is mildly \emph{harmful}: removing it
slightly improves results. The most physically evocative layer---a wave of loss
memory rippling outward---was a net negative.
\end{itemize}

Removing all four at once (the bare gradient core) costs $9.8$--$38.7\%$,
confirming that the gradient alone is not enough: it needs the scar. But it
needs \emph{only} the scar. The mass, thermostat, and ripple together account
for almost nothing, and the ripple actively hurts.

\subsection{Interpretation}

This is the parsimony result that gives the paper its title. Of the eight-layer
apparatus, the data retain exactly two mechanisms---gradient descent and the
Newton-III scar---and discard the rest. Crucially, the discarded layers were not
obviously bad ideas: effective mass, thermal scaling, and wave-like diffusion
are all physically natural, and each had a plausible story. They simply did not
survive measurement. We take this as a methodological lesson worth stating
explicitly: in physics-inspired modelling, the seductiveness of a mechanism is
no evidence of its usefulness, and only ablation tells them apart. The two-term
core of Section~\ref{sec:core} is what remained standing.
